\documentclass[12pt, titlepage]{article}
\usepackage{hyperref}

\title{Feedback Controls Lab Weekly Report \\ \normalsize{Week 11}}
\author{Aydan Vissing, Brady Schick, Gabriel Southwell, Simon Ross}
\date{\today}

\begin{document}
\maketitle

\section*{Aydan Vissing - CFO (12 hrs spent)}

I got the GUI to Milestone standard during the early week. Now that it's complete, I am going to begin work on refining it and continuing to finish the new frame design, which I have been working on since. I hope to have a print of it done by the end of the week. 

\section*{Brady Schick - CTO (10 hrs spent)}

The CMOS lab and a Poli Thought assignment took up a lot of time this week, but I was at least able to get some work done on tuning. This coming week I only have two exams to worry about, so I should have more time to work on tuning. I'm also staying on campus over the break, so I should have a lot of time to work on tuning then.

\section*{Gabriel Southwell - CSO (12 hrs spent)}

This week I worked more on integrating the tofs into the Kalman state estimator. As far as I am able to tell it should not be terribly hard to integrate. That said, it still doesn't work. I am definitely missing something in the control path so I will keep working my way through it.

On Thursday, Dr. Kohl had me install Codex in my IDE and see if it could figure it out and I am able to say that it could not. It took a lot of fussing with it to get it to actually understand what I was trying to do. It wanted to restructure everything to add the new feature. I am not sure how best to direct codex so that it is a helpful tool, but I might just end up doing it myself again. 

I hope to have xy state estimation working this week. We do need a frame made that can hold the side time of flight sensors. 

\section*{Simon Ross - CEO (12 hrs spent)}

This was a really big week for me, in both good and bad ways. On the good side of things, I got the file transfer system to set parameters from the file passed in. I am currently communicating to get this integrated into Brady's workflow. From here I just need to figure out how to use a similar system to command the drone set point.

As for the bad news, the second power regulator went out while I was working with the drone, meaning I may have some soldering to do this week. I'll communicate with Dr. Kohl on this front, and update the PCB design files accordingly. With this in mind, it's probably about time to get a drone v2 board in the mail. I'll try to get that done by the end of the week, barring any massive changes wanted by the rest of the team.

Looking forward, we've made some great progress on the app (planning to do a little more integration on that front over break), and Brady and Gabe are putting in the hours on the hard stuff. If the Lord wills, we'll have it flying soon.

\end{document}